\documentclass[12pt]{article}
\usepackage[brazil]{babel}
\usepackage[utf8]{inputenc}
\usepackage[T1]{fontenc}
\usepackage{amsmath, amssymb}
\usepackage{geometry}
\usepackage{graphicx}
\usepackage{tikz}
\usetikzlibrary{shapes,arrows,positioning,calc,matrix,arrows.meta,decorations.pathmorphing}
\usepackage{pgfplots}
\pgfplotsset{compat=1.18}
\usepackage{xcolor}
\usepackage{enumitem}
\usepackage{caption}
\usepackage{float}

\definecolor{maincolor}{RGB}{0,100,50}
\usepackage{fancyhdr}
\definecolor{forestgreen}{RGB}{0,100,50}
\usepackage[most]{tcolorbox}

\geometry{margin=2.5cm}
\pagestyle{fancy}
\fancyhf{}

\renewcommand{\headrulewidth}{2pt}
\renewcommand{\footrulewidth}{2pt}

\renewcommand{\headrule}{\hbox to\headwidth{%
  \color{black}\leaders\hrule height \headrulewidth\hfill}}

\renewcommand{\footrule}{\hbox to\headwidth{%
  \color{black}\leaders\hrule height \footrulewidth\hfill}}

\setlength{\headheight}{52.2pt} 

\fancyhead[L]{%
  \includegraphics[height=1.2cm]{Universidade_de_Rio_Verde.jpg}
}

\fancyhead[R]{%
  \raggedleft
  \textcolor{black}{%
    \small
    \textbf{Gabarito --- Grafos e Árvores}\\
    Me.~Sergio Souza Novak\\
    \textit{Lógica para Engenharia de Software}
  }
}

\fancyfoot[C]{%
  \textcolor{black}{\thepage}
}

\newcommand{\quadro}[1]{\vspace{0.2cm}\begin{tcolorbox}[colback=green!5, colframe=forestgreen, boxrule=0.5pt, arc=0pt] \textbf{Resposta:} #1 \end{tcolorbox}}

\begin{document}

\begin{enumerate}[label=\textbf{\arabic*.}]

% -------------------------------------------------------
% Questão 1
% -------------------------------------------------------
\item \textbf{Grafos Completos ($K_n$): Definição e Propriedades}
\begin{enumerate}[label=\alph*)]
    \item Desenhe o grafo $K_7$ no espaço abaixo.
    \item Qual é o grau de cada vértice em um $K_7$?
    \item Calcule o número total de arestas de um grafo completo $K_{10}$. Mostre a fórmula utilizada.
\end{enumerate}

\quadro{
\textbf{a) Desenho do $K_7$:} \\
\begin{center}
\begin{tikzpicture}[node/.style={circle, draw=black, fill=gray!20, thick, minimum size=0.6cm}]
    \foreach \i in {1,...,7} {
        \node[node] (\i) at ({360/7 * (\i - 1) + 90}:2cm) {\i};
    }
    \foreach \i in {1,...,7} {
        \foreach \j in {\i,...,7} {
            \ifnum\i<\j
                \draw[gray, thin] (\i) -- (\j);
            \fi
        }
    }
\end{tikzpicture}
\end{center}
\textbf{b)} Em um grafo completo $K_n$, todos os vértices se conectam a todos os outros. Logo, o grau de cada vértice é $n-1$. Em um $K_7$, o grau é \textbf{6}. \\
\textbf{c)} A fórmula de arestas em um grafo completo é $E = \frac{n(n-1)}{2}$. Para $n=10$: $E = \frac{10 \times 9}{2} = \frac{90}{2} =$ \textbf{45 arestas}.
}

% -------------------------------------------------------
% Questão 2
% -------------------------------------------------------
\item \textbf{Ciclos Hamiltonianos em Grafos Completos}
Utilizando a fórmula correta, calcule quantos ciclos Hamiltonianos distintos existem em um grafo completo de 6 vértices ($K_6$). Apresente a fórmula e o cálculo.

\quadro{
A fórmula para determinar o número de ciclos Hamiltonianos \textbf{distintos} (desconsiderando o sentido do percurso e o vértice inicial) em um grafo completo $K_n$ não direcionado é $\frac{(n-1)!}{2}$. \\
Para $n=6$: 
$$ \frac{(6-1)!}{2} = \frac{5!}{2} = \frac{120}{2} = \mathbf{60 \text{ ciclos Hamiltonianos distintos.}} $$
}

% -------------------------------------------------------
% Questão 5
% -------------------------------------------------------
\item \textbf{Grafos Bipartidos Completos ($K_{m,n}$)}
Desenhe o grafo bipartido completo $K_{3,2}$, separando os conjuntos em esquerda e direita.

\quadro{
\begin{center}
\begin{tikzpicture}[node/.style={circle, draw=black, fill=gray!20, thick, minimum size=0.8cm}]
    \node[node] (U1) at (0,3) {$U_1$};
    \node[node] (U2) at (0,1.5) {$U_2$};
    \node[node] (U3) at (0,0) {$U_3$};
    
    \node[node] (V1) at (4,2.25) {$V_1$};
    \node[node] (V2) at (4,0.75) {$V_2$};
    
    \draw[thick] (U1) -- (V1); \draw[thick] (U1) -- (V2);
    \draw[thick] (U2) -- (V1); \draw[thick] (U2) -- (V2);
    \draw[thick] (U3) -- (V1); \draw[thick] (U3) -- (V2);
\end{tikzpicture}
\end{center}
}

% -------------------------------------------------------
% Questão 6
% -------------------------------------------------------
\item \textbf{Aplicações: O Problema da Atribuição}
\begin{enumerate}[label=\alph*)]
    \item Modele essa situação como um grafo bipartido, desenhando os nós e as arestas.
    \item Mostre, destacando no grafo ou listando os pares, se existe um \textbf{Casamento Perfeito}...
\end{enumerate}

\quadro{
\textbf{a) Grafo Bipartido:}
\begin{center}
\begin{tikzpicture}[node/.style={rectangle, draw=black, fill=gray!20, thick, minimum width=2cm, minimum height=0.8cm, align=center}]
    % Programadores
    \node[node] (Ana) at (0,4.5) {Ana};
    \node[node] (Bob) at (0,3) {Bob};
    \node[node] (Car) at (0,1.5) {Carlos};
    \node[node] (Dan) at (0,0) {Dani};
    
    % Tarefas
    \node[node] (Fro) at (6,4.5) {Frontend};
    \node[node] (Bac) at (6,3) {Backend};
    \node[node] (IA) at (6,1.5) {IA};
    \node[node] (Dev) at (6,0) {DevOps};
    
    % Arestas originais (finas/cinza)
    \draw[gray, thick] (Ana) -- (Fro); \draw[gray, thick] (Ana) -- (IA);
    \draw[gray, thick] (Bob) -- (Bac);
    \draw[gray, thick] (Car) -- (Dev); \draw[gray, thick] (Car) -- (IA);
    \draw[gray, thick] (Dan) -- (Fro); \draw[gray, thick] (Dan) -- (Dev);
    
    % Casamento Perfeito (verdes/grossas)
    \draw[forestgreen, line width=2pt] (Bob) -- (Bac);
    \draw[forestgreen, line width=2pt] (Ana) -- (Fro);
    \draw[forestgreen, line width=2pt] (Dan) -- (Dev);
    \draw[forestgreen, line width=2pt] (Car) -- (IA);
\end{tikzpicture}
\end{center}
\textbf{b) Sim, existe um Casamento Perfeito.} \\
Bob obrigatoriamente fará \textbf{Backend}. \\
Restando IA e Frontend para Ana, e DevOps e IA para Carlos. Se Ana ficar com \textbf{Frontend}, sobra \textbf{DevOps} para Dani e \textbf{IA} para Carlos. \\
\textbf{Atribuição:} (Ana $\leftrightarrow$ Frontend), (Bob $\leftrightarrow$ Backend), (Carlos $\leftrightarrow$ IA), (Dani $\leftrightarrow$ DevOps).
}

% -------------------------------------------------------
% Questão 7
% -------------------------------------------------------
\item \textbf{Árvores Binárias: Classificação}
\begin{enumerate}[label=\alph*)]
    \item Qual a diferença fundamental entre uma Árvore Binária \textbf{Estrita} (Full) e uma Árvore Binária \textbf{Perfeita}?
    \item É possível que uma árvore seja considerada ``Completa'' mas não seja ``Perfeita''? Dê um exemplo...
\end{enumerate}

\quadro{
\textbf{a)} Em uma árvore binária \textbf{Estrita} (Full), cada nó interno deve possuir exatamente 0 ou 2 filhos (não há nós com apenas 1 filho). Já uma árvore \textbf{Perfeita} é uma árvore estrita na qual \textbf{todas as folhas estão no mesmo nível} de profundidade (formando um triângulo cheio). \\
\textbf{b) Sim}. Uma árvore completa exige apenas que todos os níveis estejam preenchidos, exceto possivelmente o último, que deve estar preenchido \textbf{da esquerda para a direita}. Abaixo, uma árvore de 4 nós: ela é completa, pois o nível 2 tem o nó mais à esquerda, mas não é perfeita (as folhas estão em níveis diferentes).
\begin{center}
\begin{tikzpicture}[level distance=0.8cm, sibling distance=1.5cm, every node/.style={circle, draw=black, thick, minimum size=0.6cm}]
    \node (1) {1}
      child {node (2) {2}
        child {node (4) {4}}
        child[missing]
      }
      child {node (3) {3}};
\end{tikzpicture}
\end{center}
}

% -------------------------------------------------------
% Questão 9
% -------------------------------------------------------
\item \textbf{Teoremas de Capacidade em Árvores}
\begin{enumerate}[label=\alph*)]
    \item Utilizando a fórmula apropriada, quantos nós \textit{no máximo} podem existir no nível $k=66$ (considerando que a raiz ocupa o nível 0)?
    \item Se uma árvore binária é perfeita e possui altura $h=312$, qual o seu número total de nós? Demonstre o cálculo.
\end{enumerate}

\quadro{
\textbf{a)} O número máximo de nós em um nível $k$ de uma árvore binária é dado por $2^k$. Portanto, no nível 66, haverá no máximo $\mathbf{2^{66}}$ nós. \\
\textbf{b)} O número total de nós em uma árvore binária \textbf{perfeita} de altura $h$ é a soma da progressão geométrica das potências de 2, cuja fórmula direta é $2^{h+1} - 1$. \\
Para $h = 312$, o número total de nós será $\mathbf{2^{313} - 1}$.
}

% -------------------------------------------------------
% Questão 10
% -------------------------------------------------------
\item \textbf{Pior Caso (Degeneração) e a Relação das Arestas}
\begin{enumerate}[label=\alph*)]
    \item Desenhe uma árvore degenerada com 5 nós. Por que cientistas da computação consideram essa formação um ``pior caso'' para a busca (qual a sua notação Big-O)?
    \item Baseando-se no Lema do Aperto de Mãos aplicado a árvores (Grafos Acíclicos Conexos), quantas arestas obrigatoriamente existem em uma árvore binária que possui exatos 150 vértices? Justifique matematicamente porque não pode haver uma aresta a mais nem uma a menos.
\end{enumerate}

\quadro{
\textbf{a)} Uma árvore degenerada é semelhante a uma lista encadeada simples:
\begin{center}
\begin{tikzpicture}[node/.style={circle, draw=black, thick, minimum size=0.6cm}]
    \node[node] (1) at (0,0) {1};
    \node[node] (2) at (1.5,0) {2};
    \node[node] (3) at (3,0) {3};
    \node[node] (4) at (4.5,0) {4};
    \node[node] (5) at (6,0) {5};
    \draw[->, thick] (1) -- (2); \draw[->, thick] (2) -- (3);
    \draw[->, thick] (3) -- (4); \draw[->, thick] (4) -- (5);
\end{tikzpicture}
\end{center}
É considerada o pior caso pois a árvore perde seu formato ramificado (que daria complexidade $O(\log n)$) e as buscas passam a percorrer a estrutura linearmente, exigindo \textbf{tempo linear $O(n)$}. \\
\textbf{b)} Em toda árvore (grafo conexo e sem ciclos), o número de arestas $E$ sempre obedece rigorosamente à relação $E = V - 1$. Portanto, para 150 vértices, há exatamente \textbf{149 arestas}. \\
\textbf{Justificativa:} Se houver uma aresta a menos (148 arestas), o grafo fatalmente se desconecta e deixa de ser árvore (ficará com componentes separadas). Se houver uma aresta a mais (150 arestas), fatalmente se formará um ciclo, o que também contraria a definição de árvore.
}

\end{enumerate}

\end{document}
